\subsection{Elementary criteria for self-adjoint operators}

PROPOSITION 9. --- Let \(v \in \mathcal{L}(\mathrm{F};\mathrm{E})\) be a continuous injective linear map from F into E whose image is dense in E. The adjoint of \(v\) is a continuous injective linear map from E into F and one has \((v^{*})^{-1} = (v^{-1})^{*}\) . In particular, if E = F, the endomorphism \(v\) is Hermitian if and only if the partial operator \(v^{-1}\) is self-adjoint.

The partial operator \(v^{-1}\) is a closed operator with dense domain from E into F (example 2 of IV, p. 228) and the adjoint \(v^{*}\) of v is a continuous linear map from E into F; it is injective, since the image of v is dense in E (EVT, V, p. 41, prop. 4). Let s (resp. \(s'\) ) be the isometric isomorphism \((x,y)\mapsto(-y,x)\) from E ⊕ F onto F ⊕ E (resp. the isometric isomorphism \((y,x)\mapsto(-x,y)\) from F ⊕ E onto E ⊕ F) and let \(\iota\) (resp. \(\iota'\) ) be the isometric isomorphism \((y,x)\mapsto(x,y)\) from F ⊕ E onto E ⊕ F (resp. the isometric isomorphism \((x,y)\mapsto(y,x)\) from E ⊕ F onto F ⊕ E). One then has \(s\circ\iota=-\iota'\circ s'\) , whence

\[
\Gamma_ {(v ^ {- 1}) ^ {*}} = s \left(\Gamma_ {v ^ {- 1}}\right) ^ {\circ} = s \left(\iota \left(\Gamma_ {v}\right)\right) ^ {\circ} = - \iota^ {\prime} \left(s ^ {\prime} \left(\Gamma_ {v}\right)\right) ^ {\circ} = - \iota^ {\prime} \left(\Gamma_ {v ^ {*}}\right) = \Gamma_ {(v ^ {*}) ^ {- 1}}
\]

by prop. 7 of IV, p. 236. The proposition follows.

PROPOSITION 10 (Hellinger--Toeplitz). --- Let u be a symmetric partial operator on the Hilbert space E. If the domain of u is equal to E, then \(u \in \mathcal{L}(E)\) and u is Hermitian.

Indeed, the partial operator u is closable (prop. 8 of IV, p. 237), and hence closed since its domain is E (IV, p. 228, remark). One then concludes by invoking EVT, I, p. 19, cor. 5.

COROLLARY. --- Let u be a symmetric partial operator on E. If the partial operator u induces a bijective linear map from dom(u) onto E, then u is self-adjoint.

Indeed, the partial operator \(u^{-1}\) inverse to u is symmetric with domain E (prop. 9), hence \(u^{-1}\) is a self-adjoint element of \(\mathcal{L}(\mathrm{E})\) (prop. 10), and u is then Hermitian (prop. 9).

PROPOSITION 11. --- Let u be a symmetric partial operator on E and let \(\lambda \in \mathbf{C}\) . If \(u + \lambda 1_{\mathrm{E}}\) and \(u + \overline{\lambda} 1_{\mathrm{E}}\) are surjective, then u is self-adjoint.

It suffices to show that \(\operatorname{dom}(u^{*}) \subset \operatorname{dom}(u)\) Let \(x \in \operatorname{dom}(u^{*})\) . By hypothesis there exists \(y \in \operatorname{dom}(u)\) such that \(u(y) + \overline{\lambda}y = u^{*}(x) + \overline{\lambda}x\) . Let us show that y = x. For all \(z \in \operatorname{dom}(u)\) , it follows that

\[
\begin{array}{r l} & {\langle (u + \lambda 1 _ {\mathrm{E}}) (z) | x \rangle = \langle z | (u ^ {*} + \overline {{\lambda}} 1 _ {\mathrm{E}}) (x) \rangle} \\ & {\qquad = \langle z | (u + \overline {{\lambda}} 1 _ {\mathrm{E}}) (y) \rangle = \langle (u + \lambda 1 _ {\mathrm{E}}) (z) | y \rangle ,} \end{array}
\]

since u is symmetric. Since the operator \(u + \lambda1_{E}\) is surjective, one indeed has y = x, hence \(x \in \text{dom}(u)\) .

We shall see later (cf.prop. 17 of IV, p. 248) that if \(\lambda \in \mathbf{C} - \mathbf{R}\) , then the converse is true.

PROPOSITION 12. --- Let u be a closed operator with dense domain from E into F. The partial operator \(u^{*}\circ u\) on E is self-adjoint and positive. Its domain is a core for u.

Denote by \(v\) the partial operator the partial operator \(1_{\mathrm{E}} + u^{*} \circ u\) . Its domain is \(\operatorname{dom}(u^{*} \circ u)\) , which is contained in \(\operatorname{dom}(u)\) . For all \(x \in \operatorname{dom}(u)\) and \(y \in \operatorname{dom}(v)\) , we have \(u(y) \in \operatorname{dom}(u^{*})\) and

\begin{enumerate}
\def\labelenumi{(\arabic{enumi})}
\setcounter{enumi}{2}
\tightlist
\item
\end{enumerate}

\[
\langle x \mid v (y) \rangle = \langle x \mid y \rangle + \langle x \mid (u ^ {*} \circ u) (y) \rangle = \langle x \mid y \rangle + \langle u (x) \mid u (y) \rangle ,\tag{4}
\]

\[
\langle y \mid v (y) \rangle = \| y \| ^ {2} + \| u (y) \| ^ {2}.
\]

Formula (4) implies that the partial operator \(v\) is injective. Moreover, by prop. 7 of IV, p. 236, \(a\) ), one has \(\mathrm{F} \oplus \mathrm{E} = \Gamma_{u^*} \oplus s(\Gamma_u)\) Let \(x \in \mathrm{E}\) . There exists \(y' \in \mathrm{dom}(u^*)\) and \(x' \in \mathrm{dom}(u)\) such that

\[
(0, x) = (y ^ {\prime}, u ^ {*} (y ^ {\prime})) + (- u (x ^ {\prime}), x ^ {\prime}) = (y ^ {\prime} - u (x ^ {\prime}), x ^ {\prime} + u ^ {*} (y ^ {\prime})).
\]

We therefore have \(y' = u(x')\) , whence \(x' \in \text{dom}(u^* \circ u) = \text{dom}(v)\) and

\[
x = x ^ {\prime} + u ^ {*} (y ^ {\prime}) = x ^ {\prime} + (u ^ {*} \circ u) (x ^ {\prime}) = v (x ^ {\prime}).
\]

The partial operator v on E is therefore surjective, and induces a bijective linear map from dom(v) to E.

Let \(x \in \mathrm{E}\) orthogonal to \(\operatorname{dom}(v)\) . Let us write \(x = v(x')\) where \(x' \in \operatorname{dom}(v)\) . By formula (4), we obtain

\[
0 = \langle x ^ {\prime} | x \rangle = \langle x ^ {\prime} | v (x ^ {\prime}) \rangle = \| x ^ {\prime} \| ^ {2} + \| u (x ^ {\prime}) \| ^ {2}
\]

whence \(x' = 0\) , then \(x = 0\) . The domain of \(v\) is therefore dense in E.

The partial operator \(v\) has dense domain; it is bijective and the formula (3) shows that it is symmetric. We conclude that \(v\) is self-adjoint by applying the corollary of proposition 10. Consequently, \(u^{*} \circ u = v - 1_{\mathrm{E}}\) is self-adjoint. Moreover, the formula

\[
\langle x \mid (u ^ {*} \circ u) (x) \rangle = \| u (x) \| ^ {2}
\]

for all \(x \in \text{dom}(u^{*} \circ u)\) implies that \(u^{*} \circ u\) is positive.

Finally, let \(y \in \mathrm{E}_u\) orthogonal to \(\operatorname{dom}(u^* \circ u)\) . There exists an element \(x\) in the domain of \(u^* \circ u\) such that \(y = v(x) = x + (u^* \circ u)(x)\) . We then have \(0 = (x \mid y)_u = \langle y \mid y \rangle\) , whence \(y = 0\) .

